\documentclass[11pt, headheight=30pt]{scrartcl}
\usepackage[white, nologo]{darkWhite}

\title{18.675 Problem Set 1}

\begin{document}
% \maketitle
\section{$\sigma$-algebra manipulation}
\begin{prob}
    [Unions and Intersections]
    Let $E$ be a set and let $S$ be a set of $\sigma$-algebras on $E$. Define
    \[
        \mathcal{E}^* = \{A \subseteq E : A \in \mathcal{E} \text{ for all } \mathcal{E} \in S\}.
    \]
    Show that $\mathcal{E}^*$ is a $\sigma$-algebra on $E$.
    Show, on the other hand, by example, that the union
    of two $\sigma$-algebras on the same set need not be a $\sigma$-algebra.
\end{prob}

Just check the axioms:
\begin{enumerate}[(a)]
    \item $\emptyset \in \mathcal{E}$ for all $\mathcal{E} \in S$,
    so $\emptyset \in \mathcal{E}^*$.
    \item If $A \in \mathcal{E}^*$, then $A \in \mathcal{E}$
    for all $\mathcal{E} \in S$.
        But then $A^c \in \mathcal{E}$ for all $\mathcal{E} \in S$,
        so $A^c \in \mathcal{E}^*$.
    \item If $(A_n : n \in \NN)$ is a sequence of sets in $\mathcal{E}^*$, then
        $A_n \in \mathcal{E}$ for all $\mathcal{E} \in S$ and all $n \in \NN$.
        But then $\bigcup_n A_n \in \mathcal{E}$ for all $\mathcal{E} \in S$,
        so $\bigcup_n A_n \in \mathcal{E}^*$.
\end{enumerate}

The counterexample is really stupid: let $E = \{1,2,3,4\}$, let
$\mathcal{E}_1 = \{\emptyset, \{1,2\}, \{3,4\}, E\}$, and let
$\mathcal{E}_2 = \{\emptyset, \{1,3\}, \{2,4\}, E\}$. Then
$\mathcal{E}_1 \cup \mathcal{E}_2 = \{\emptyset, \{1,2\}, \{1,3\}, \{2,4\}, \{3,4\}, E\}$
is not closed under countable unions (where is $\{1,2,3\}$?).

\section{Real number computations}
\begin{prob}
    Show that the following sets of subsets of $\RR$ all generate
    the same $\sigma$-algebra:
    \begin{enumerate}[(a)]
        \item $\{(a, b) : a < b\}$,
        \item $\{(a, b] : a < b\}$,
        \item $\{(-\infty, b] : b \in \RR\}$.
    \end{enumerate}
\end{prob}

For notation, let
$\mathcal{E}_a, \mathcal{E}_b, \mathcal{E}_c$ denote the $\sigma$-algebras
generated by $(a), (b), (c)$ respectively.

\begin{psec*}
    [Proof $\mathcal{E}_c\subset \mathcal{E}_a$]
    By complements, $\mathcal{E}_a$ contains all sets of the form
    \[
        (a,b]^c = (-\infty, a] \cup [b, \infty)
    \]
    for all $a < b$. By countable intersections, we have
    \[
        \bigcup_{b\in \ZZ, b > a} (-\infty, a] \cup [b, \infty) = (-\infty, a]  
    \]
    thus $\mathcal{E}_c\subset \mathcal{E}_a$.
\end{psec*}

\begin{psec*}
    [Proof $\mathcal{E}_a\subset \mathcal{E}_b$]
    By countable unions, $\mathcal{E}_b$ contains all sets of the form
    \[
        \bigcup_{n\in \NN} (a, b - 1/n] = (a,b)
    \]
    for all $a < b$. Thus $\mathcal{E}_a\subset \mathcal{E}_b$.
\end{psec*}

\begin{psec*}
    [Proof $\mathcal{E}_b\subset \mathcal{E}_c$]
    By complements and intersections, $\mathcal{E}_c$ contains all sets of the form
    \[
        (-\infty, b] \setminus (-\infty, a] = (a,b]
    \]
    for all $a < b$. Thus $\mathcal{E}_b\subset \mathcal{E}_c$.
\end{psec*}


\begin{prob}
    Show that a countably additive set function on a ring $\mathcal{S}$
    is additive, increasing and countably subadditive.
\end{prob}

Countably additive sets are automatically
additive. This is also clearly increasing because
for any $A, B\in \mathcal{S}$ such that $A\subset B$,
$B\setminus A\in \mathcal{S}$, thus
\[
    \mu\left(B\right) = \mu\left(A\right) + \mu\left(B\setminus A\right)
    \geq \mu\left(A\right).    
\]

Finally, we show countable subadditivity. For all sequences
$A_1, A_2, \cdots \in \mathcal{S}$
such that $\bigcup_n A_n \in \mathcal{S}$, we know that
$B_i = A_i \setminus \bigcup_{n < i} A_n \in \mathcal{S}$
for all $i\geq 1$, and of course
\[
    \bigcup_n A_n = \bigcup_n B_n.    
\]
On the other hand, the $B_n$ are disjoint by construction, so
\[
    \mu\left(\bigcup_n A_n\right) = \mu\left(\bigcup_n B_n\right)
    = \sum_n \mu\left(B_n\right) \leq \sum_n \mu\left(A_n\right).
\]
where the last statement is because $B_n\subset A_n$
and we know $\mu$ is increasing.

\section{Dynkin systems}

\begin{prob}
    Show that a $\pi$-system which is also a $d$-system is a $\sigma$-algebra.
\end{prob}

First off, all $\sigma$-algebras are clearly $\pi$ systems.
They are $d$-systems, because:
\begin{enumerate}[(a)]
    \item $\Omega$ is in the $\sigma$-algebra.
    \item If $A\subset B$ are in the $\sigma$-algebra, then
    $B\setminus A = B\cap A^c$ is also in the $\sigma$-algebra.
    \item If $A_1, A_2, \cdots$ are in the $\sigma$-algebra, then
    $\bigcup_n A_n$ is also in the $\sigma$-algebra
    (This is stronger than we need; a $d$-system needs only be closed under
    countable increasing union).
\end{enumerate}

If you are a $d$ system, then you have $\Omega$.
Thus by the existence of proper differences you have complements. If you
also a $\pi$ system, then finite intersections and complements give you
finite unions via De Morgan's law.

Finally, given any collection $A_1, A_2,\cdots \in \LL$,
their prefix unions $A_1\cup \cdots A_n$ are all in $\LL$, and thus by
the countable increasing union condition of $d$-systems,
their union is also in $\LL$.

\begin{prob}
    [Countable Additivity Condition]
    Let $\mu$ be a finite-valued additive set function on a ring $\mathcal{A}$.
    Show that $\mu$ is countably additive if and only if the following
    condition holds: for any decreasing sequence $(A_n : n \in \NN)$
    of sets in $\mathcal{A}$, with $\bigcap_n A_n = \emptyset$, we have
    $\mu(A_n) \to 0$ as $n \to \infty$.
\end{prob}

\begin{psec*}
    [Proof $\Rightarrow$]
    Suppose $\mu$ is countably additive. In particular we're additive,
    so we can let $B_i = A_{i}\setminus A_{i+1}\in \mathcal{A}$ for all $i\geq 1$;
    these are all disjoint by construction.
    Then by countable additivity, we know that
    \[
        \infty > \mu(A_n) = \mu\left(\bigcup_{i\geq n} B_i\right) = \sum_{i\geq n} \mu(B_i)  
    \]
    for all $n\geq 1$. By definition of infinite sums, for all
    $\eps > 0$, there exists an $M$ such that
    \[
        \sum_{n\leq i\leq n + M} \mu(B_i) > \sum_{i\geq n} \mu(B_i) - \eps \implies \sum_{n + M\leq i} \mu(B_i) < \eps.
    \]
    Thus $\mu(A_n) \to 0$ as $n\to \infty$.
\end{psec*}

\begin{psec*}
    [Proof $\Leftarrow$]
    Suppose we know the second condition. Then, for any collection
    of disjoint sets $A_1, A_2, \cdots \in \mathcal{A}$ such
    that $A = \bigcup_n A_n \in \mathcal{A}$, we can define, we can define
    \[
        B_n = \bigcup_{i\geq n} A_i = A\setminus\left( \bigcup_{i < n} A_i \right) \in \mathcal{A}
    \]
    Clearly, $B_n$ form a decreasing sequence such that
    $\bigcap_n B_n = \emptyset$, thus $\mu(B_n) \to 0$ as $n\to \infty$.

    But this is what we want. First off, for any $\eps$, there exists an $N$
    such that $\mu(B_N) < \eps$, thus
    \[
        \eps + \sum_{i\in \NN} \mu(A_n) \geq \mu(B_N) + \sum_{i< N} \mu(A_n) = \mu(A)
    \]
    thus $\sum_{i\in \NN} \mu(A_n) \geq \mu(A)$. On the other hand,
    for all $N$,
    \[
        \sum_{i\leq N} \mu(A_i) = \mu\left(\bigcup_{i\leq N} A_i\right) \leq \mu(A)    
    \]
    (the inequality is because additive set functions
    on rings are automatically increasing), thus we can pass to the limit.
\end{psec*}

\section{Interesting things}

\begin{prob}[Baby Fatou-Lebesgue Theorem]
    Let $(E, \mathcal{E}, \mu)$ be a finite measure space. Show that, for
    any sequence of sets $(A_n : n \in \NN)$ in $\mathcal{E}$,
    \[
        \mu\left(\liminf_n A_n\right) \leq \liminf_n \mu(A_n) \leq
        \limsup_n \mu(A_n) \leq \mu\left(\limsup_n A_n\right).
    \]
    Show that the first inequality remains true without the
    assumption that $\mu(E) < \infty$, but that the last
    inequality may then be false.
\end{prob}
The middle inequality is true because for any set $S$,
$\inf S \leq \sup S$.

\begin{claim*}
    [Continuity from Above and Below]
    We write $A_i \uparrow A$ if $A_1 \subseteq A_2 \subseteq \dots$
    and $A = \bigcup_i A_i$. Then,
    \[
        \mu(A) = \lim_i \mu(A_i).
    \]

    Similarly, we write $A_i \downarrow A$ if $A_1 \supseteq A_2 \supseteq \dots$
    and $A = \bigcap_i A_i$. If $\mu(A_1) < \infty$, then
    \[
        \mu(A) = \lim_i \mu(A_i).
    \]
\end{claim*}
\begin{proof}
    Let $A = \bigcup_i (A_i \setminus A_{i-1})$, each of which is disjoint.
    Then
    \[
        \mu(A) = \sum_i \mu(A_i \setminus A_{i-1})
        = \lim_i \sum_{j\leq i} \mu(A_j \setminus A_{j-1}) = \lim_i \mu(A_i).    
    \]

    For the $\downarrow$ direction, apply the previous
    result to $A_1 \setminus A_i$ to yield
    \[
        \mu(A_1\setminus A) = \lim_i \mu(A_1 \setminus A_i)
        = \lim_i \mu(A_1) - \mu(A_i) = \mu(A_1) - \lim_i \mu(A_i)
    \]
    We need $\infty > \mu(A_1) > \mu(A_i)$ for all the subtractions
    to avoid $\infty - \infty$.
\end{proof}

    
\begin{psec*}
    [First inequality]
    We do not assume $\mu(E) < \infty$. First we construct
    \[
        B_n =  \bigcap_{k\geq n} A_k \in \mathcal{E}   
    \]
    (note that these are increasing; this is the usual $\sup \inf = \lim \inf$
    identity).
    % Thus let
    % \[
    %     B = \bigcup_{n\geq 1} B_n \in \mathcal{E}    
    % \]
    Now the claim applies, so we must show
    \[
        \lim_{n\to \infty} \mu\left(B_n\right) \leq \lim_{n\to \infty} \inf_{k\geq n} \mu(A_k)    
    \]
    Actually, this is just true termwise:
    \[
        \mu\left(\bigcap_{k\geq n} A_k\right) \leq \mu\left(A_i\right)
    \]
    for all $i \geq n$, thus
    \[
        \mu\left(\bigcap_{k\geq n} A_k\right) \leq \inf_{k\geq n} \mu(A_k)
    \]
\end{psec*}


\begin{psec*}
    [Third inequality]
    Do the same thing.
    \[
        B_n =  \bigcup_{k\geq n} A_k \in \mathcal{E}   
    \]
    These are decreasing and $\mu(B_1) < \mu(E) < \infty$,
    so we apply the claim: we want
    \[
        \lim_{n\to \infty} \mu\left(B_n\right) \geq \lim_{n\to \infty} \sup_{k\geq n} \mu(A_k)    
    \]
    This is again true termwise:
    \[
        \mu\left(\bigcup_{k\geq n} A_k\right) \geq \mu\left(A_i\right)
    \]
    for all $i \geq n$, thus
    \[
        \mu\left(\bigcup_{k\geq n} A_k\right) \geq \sup_{k\geq n} \mu(A_k)
    \]
\end{psec*}

\begin{prob}
    Let $(A_n : n \in \NN)$ be a sequence of events in a
    probability space $\Omega$. Show that the events $A_n$
    are independent if and only if the $\sigma$-algebras
    $\sigma(A_n) = \{\emptyset, A_n, A_n^c, \Omega\}$ are independent.
\end{prob}

If the $\sigma$-algebras
are independent, then by definition for any $J\subset \NN$,
\[
    \mu\left(\bigcap_{i\in J} A_i\right) = \prod_{i\in J} \mu(A_i)
\]
thus the events are independent.

We now want a stupid claim: if $(A_n : n\leq N)$ are a finite
sequence of independent events, then $\{A_1^c, A_2, A_3,\cdots A_n\}$
(where we have $A_1^c$ instead of $A_1$) are also independent. This is
clear; if $J \subset \{2,\dots, N\}$, then nothing has changed;
if $J = \{1\}\cup K$ with $K\subset \{2,\dots, N\}$, then
\begin{align*}
    \mu\left(A_1^c\cap \bigcap_{i\in K} A_i\right) &= \mu\left(\bigcap_{i\in K} A_i\right) - \mu\left(A_1\cap \bigcap_{i\in K} A_i\right)\\
    &= \prod_{i\in K} \mu(A_i) - \mu(A_1)\cdot \prod_{i\in K} \mu(A_i)\\
    &= \mu(A_1^c)\cdot \prod_{i\in K} \mu(A_i)    
\end{align*}
as desired. By repeated applying this a finite number of
times, we find that any finite subcollection of the $A_n$
can be replaced by its complements without changing
the independence.

Given the events are independent, lets consider any family of
sets $B_n \in \sigma(A_n)$ for $n\in \NN$. We want to show
\[
    \mu\left(\bigcap_{n\in \NN} B_n\right) = \prod_{n \in \NN} \mu(B_n)
\]
By continuity from above and definition of infinite product,
it suffices to show
\[
    \mu\left(\bigcap_{n\leq N} B_n\right) = \prod_{n\leq N} \mu(B_n)    
\]
for all $N$. We induct; the base case $n = 0$ is clear.

Then, assume we know the $n-1$ case. If any $B_i$ is $\emptyset$,
both sides are $0$. If any $B_i$ is $\Omega$, we can remove
it from the expression without changing either side, and
reduce to the $n-1$ case. Thus we can assume all $B_i$ are
nontrivial. By the claim, we are now done.

\begin{prob}
    Let $B$ be a Borel subset of the interval $[0, 1]$.
    Show that for every $\epsilon > 0$, there exists a finite
    union of disjoint intervals $A = (a_1, b_1] \cup \cdots \cup (a_n, b_n]$
    such that the Lebesgue measure of $A \triangle B = (A^c \cap B) \cup (A \cap B^c)$
    is less than $\epsilon$.
    Show further that this remains true for every Borel set in $\RR$
    of finite Lebesgue measure.
\end{prob}

Let $\mathcal{P} \subset \mathcal{B}(\RR)$ contain all Borel subsets $B$ satisfying
the claim. We claim that $\mathcal{P}$ is a $\pi$-system. Indeed, for any
$C,D\in \mathcal{P}$, for every $\eps > 0$, there exists $A$ and $B$ which
are finite unions of disjoint intervals such that
\[
    \mu\left(A\triangle C\right), \mu\left(B\triangle D\right) < 0.0001\eps
\]

Then, $A\cap B$ is a finite union of disjoint intervals.\footnote{
    Okay we can verify this; if $A = E_1\cup \cdots \cup E_n$ and
    $B = F_1 \cup \cdots \cup F_m$ are finite unions of disjoint intervals,
    then $A\cap B = E_1\cap F_1 \cup \cdots \cup E_n\cap F_m$ is a finite
    union of disjoint intervals, some of which might be empty.
} We now bash some inequalities:
\begin{align*}
    \mu\left\{(A\cap B)\triangle (C\cap D)\right\} &= \mu\left\{\left((A\cap B)\cap (C^c\cup D^c)\right) \cup \left((C\cap D)\cap (A^c\cup B^c)\right)\right\}\\
    &= \mu\left\{\left((A\cap B)\cap C^c\right) \cup  \left((A\cap B)\cap D^c\right)\right.\\
    &\left.\cup \left((C\cap D)\cap A^c\right) \cup \left((C\cap D)\cap B^c\right)\right\}\\
    &\leq \mu\left((A\cap B)\cap C^c\right) + \mu\left((A\cap B)\cap D^c\right)\\
    &+ \mu\left((C\cap D)\cap A^c\right) + \mu\left((C\cap D)\cap B^c\right)\\
    &\leq \mu\left(A\cap C^c\right) + \mu\left(B\cap D^d\right) + \mu\left(C\cap A^c\right) + \mu\left(D\cap B^c\right)\\
    &< 0.0002\eps
\end{align*}

\section{Caratheodory yields complete measures}

\begin{prob}
    Let $(E, \mathcal{E}, \mu)$ be a measure space. Call a subset $N \subseteq E$ null if
    $N \subseteq B$ for some $B \in \mathcal{E}$ with $\mu(B) = 0$.
    Write $\mathcal{N}$ for the set of null sets.
    Prove that the set of subsets
    \[
        \mathcal{E}_\mu = \{A \cup N : A \in \mathcal{E}, N \in \mathcal{N}\}
    \]
    is a $\sigma$-algebra and show that $\mu$ has a well-defined and countably
    additive extension to $\mathcal{E}_\mu$ given by $\mu(A \cup N) = \mu(A)$.
    We call $\mathcal{E}_\mu$ the completion of $\mathcal{E}$ with respect to
    $\mu$.
    
    Suppose now that $\mathcal{E}$ is $\sigma$-finite and write $\mu^*$
    for the outer measure associated to $\mu$, as in the proof of Caratheodory's
    extension theorem. Show that $\mathcal{E}_\mu$ is exactly the
    set of $\mu^*$-measurable sets.
\end{prob}

\begin{psec*}
    [Proof $\mathcal{E}_\mu$ is a $\sigma$-algebra]
    We check the axioms:
    \begin{enumerate}[(a)]
        \item $\emptyset \in \mathcal{E}$ and $\emptyset \in \mathcal{N}$,
        so $\emptyset \in \mathcal{E}_\mu$.
        \item In the above definition where $N\subset B \in \mathcal{E}$ with $\mu(B) = 0$,
        note that $B\setminus N \subset B$ too. This lets us write
        \[
            \left(A\cup N\right)^c = A^c \setminus N
            = \underbrace{A^c\setminus B}_{\in \mathcal{E}} \cap \underbrace{B\setminus N}_{\in \mathcal{N}}
        \]
        thus $\left(A\cup N\right)^c \in \mathcal{E}_\mu$.

        \item If $(A_n\cup N_n : n\in \NN)$ is a sequence of sets in $\mathcal{E}_\mu$,
        with each $N_n\subset B_n$ with $\mu(B_n) = 0$,
        then
        \[
            \bigcup_n \left(A_n\cup N_n\right) = \underbrace{\bigcup_n A_n}_{\in \mathcal{E}} \cup \underbrace{\bigcup_n N_n}_{\subset \bigcup_n B_n}
        \]
        Note that $\mu\left(\bigcup_n B_n\right)\leq \sum_n \mu(B_n) = 0$,
        so $\bigcup_n \left(A_n\cup N_n\right) \in \mathcal{E}_\mu$.
    \end{enumerate}
\end{psec*}

\begin{psec*}
    [Extension]
    To show this is well-defined, we need to show that the relation $A\sim C$
    iff there exists $N, M\in \mathcal{N}$ such that $A\cup N = C\cup M$ is an equivalence
    relation. Reflexivity and symmetry are clear. For transitivity, if
    $A\cup N = C\cup M$ and $C\cup P = D\cup Q$, then
    \[
        A\cup \left(N \cup P\right) = C\cup M \cup P = D\cup \left(Q \cup M\right)
    \]
    where $N\cup P$ and $Q\cup M$ are both null, so $A\sim D$.
\end{psec*}

\begin{psec*}
    [Countable Additivity]
    Let $(A_n\cup N_n : n\in \NN)$ be a sequence of disjoint sets in $\mathcal{E}_\mu$.
    Then
    \begin{align*}
        \mu\left(\bigcup_n \left(A_n\cup N_n\right)\right) &= \mu\left(\bigcup_n A_n \cup \bigcup_n N_n\right)\\
        &= \mu\left(\bigcup_n A_n\right)\\
        &= \sum_{n\in \NN} \mu(A_n)\\
        &= \sum_{n\in \NN} \mu\left(A_n\cup N_n\right)
    \end{align*}
\end{psec*}

\todo{Finish this!}


\section{Function manipulation}

\begin{prob}[Preliminaries]
    Let $(f_n: n \in \NN)$ be a sequence of measurable functions on a measurable
    space $(E, \mathcal{E})$. Show that the following functions are also measurable:
    \begin{enumerate}[(a)]
        \item $f_1 + f_2$,
        \item $f_1 f_2$,
        \item $\inf_n f_n$,
        \item $\sup_n f_n$,
        \item $\liminf_n f_n$,
        \item $\limsup_n f_n$.
    \end{enumerate}
    Show also that $\{x \in E : f_n(x) \text{ converges as } n \to \infty\} \in \mathcal{E}$.
\end{prob}


\begin{prob}
    [Pushforward measure]
    Let $(E, \mathcal{E})$ and $(G, \mathcal{G})$ be measurable spaces,
    let $\mu$ be a measure on $E$, and let $f: E \to G$ be a measurable
    function. Show that we can define a measure $\nu$ on $G$ by setting
    $\nu(A) = \mu(f^{-1}(A))$ for each $A \in \mathcal{G}$.
\end{prob}

\section{Independence}

\begin{prob}
    Show that the following condition implies that the random variables $X$ and $Y$
    are independent: $\PP(X \leq x, Y\leq y) = \PP(X\leq x)\PP(Y\leq y)$
    for all $x, y \in \RR$.
\end{prob}

\end{document}